% !TeX program = lualatex
% !TeX root = main.tex


\documentclass[12pt,a4paper]{article}

\usepackage{fontspec}
\usepackage{polyglossia}
\setdefaultlanguage{russian}
\setotherlanguage{english}

\setmainfont{CMU Serif}

\usepackage{amsmath}
\usepackage{amssymb}
\usepackage{geometry}

\geometry{margin=25mm}

\title{
    Физические основы численного
    моделирования электромагнитных волн
}

\author{Кирилл Кабаев}
\date{\today}

\begin{document}

\maketitle

\begin{abstract}
Рассматриваются физические основы
вычислительной электродинамики,
уравнения Максвелла и численные методы
моделирования волновых процессов.
\end{abstract}

\section{Уравнения Максвелла}

Закон электромагнитной индукции Фарадея:

\begin{equation}
\nabla \times \mathbf{E}
=
-\frac{\partial \mathbf{B}}{\partial t}.
\end{equation}

\section{Уравнение Гельмгольца}

В однородной среде для скалярного
волнового поля:

\begin{equation}
\nabla^2 u + k^2 u = 0.
\end{equation}

\end{document}
